% GENERATED FILE. Edit canonical lessons, then run npm run generate:books.
% canonical-source-hash: fnv1a64:798a3ea25ed6a92a
% canonical-lessons: PT-C242-cozido-cozida, PT-C242-maduro-madura, PT-C242-escuro-escura, PT-C242-claro-clara, PT-C242-distante

\chapter{Cooked, Ripe, Dark, Clear, Distant}
\label{ch:pt-cooked-ripe-dark-clear-distant}
\hlchaptermodality{242}

\begin{chapteropening}
\textbf{By the end of this chapter:} \emph{I can say cooked, ripe, dark, clear and distant.}

Complete the last lesson of chapter 242: I can say cooked, ripe, dark, clear and distant.
\end{chapteropening}

\section[cozido / cozida]{cozido / cozida — cooked, boiled}

\label{lesson:PT-C242-cozido-cozida}



\begin{quote}
Before the new one: say the Portuguese for unknown, then the Portuguese for simple.
\end{quote}

\subsection*{You'll want to know: cozido / cozida}
\textbf{cozido / cozida} — \textquotedblleft{}cooked, boiled\textquotedblright{}.

\begin{grammarlens}[title={What you've built}]
One more word for this chapter.
\end{grammarlens}

\subsection*{Your turn}
\begin{itemize}
\raggedright
  \item \emph{Say it:} \emph{cozido / cozida}
  \item \emph{Say it:} \emph{cozido / cozida}, once more
  \item \emph{Say it:} say \emph{cozido / cozida}
  \item \emph{Recall:} say the Portuguese for unknown, then the Portuguese for simple, then say \emph{cozido / cozida} again
\end{itemize}

\begin{tcolorbox}[breakable,colback=teal!4,colframe=teal!35!black,title={Before you move on}]
What does cozido / cozida mean? (\textquotedblleft{}Cooked, boiled\textquotedblright{}.) Say it once more.
\end{tcolorbox}
\section[maduro / madura]{maduro / madura — ripe; mature}

\label{lesson:PT-C242-maduro-madura}



\begin{quote}
Before the new one: say the Portuguese for simple, then the Portuguese for cooked.
\end{quote}

\subsection*{You'll want to know: maduro / madura}
\textbf{maduro / madura} — \textquotedblleft{}ripe; mature\textquotedblright{}.

\begin{grammarlens}[title={What you've built}]
One more word for this chapter.
\end{grammarlens}

\subsection*{Your turn}
\begin{itemize}
\raggedright
  \item \emph{Say it:} \emph{maduro / madura}
  \item \emph{Say it:} \emph{maduro / madura}, once more
  \item \emph{Say it:} say \emph{maduro / madura}
  \item \emph{Recall:} say the Portuguese for simple, then the Portuguese for cooked, then say \emph{maduro / madura} again
\end{itemize}

\begin{tcolorbox}[breakable,colback=teal!4,colframe=teal!35!black,title={Before you move on}]
What does maduro / madura mean? (\textquotedblleft{}Ripe; mature\textquotedblright{}.) Say it once more.
\end{tcolorbox}
\section[escuro / escura]{escuro / escura — dark}

\label{lesson:PT-C242-escuro-escura}



\begin{quote}
Before the new one: say the Portuguese for cooked, then the Portuguese for ripe; mature.
\end{quote}

\subsection*{You'll want to know: escuro / escura}
\textbf{escuro / escura} — \textquotedblleft{}dark\textquotedblright{}.

\begin{grammarlens}[title={What you've built}]
One more word for this chapter.
\end{grammarlens}

\subsection*{Your turn}
\begin{itemize}
\raggedright
  \item \emph{Say it:} \emph{escuro / escura}
  \item \emph{Say it:} \emph{escuro / escura}, once more
  \item \emph{Say it:} say \emph{escuro / escura}
  \item \emph{Recall:} say the Portuguese for cooked, then the Portuguese for ripe; mature, then say \emph{escuro / escura} again
\end{itemize}

\begin{tcolorbox}[breakable,colback=teal!4,colframe=teal!35!black,title={Before you move on}]
What does escuro / escura mean? (\textquotedblleft{}Dark\textquotedblright{}.) Say it once more.
\end{tcolorbox}
\section[claro / clara]{claro / clara — clear, pale (of a colour)}

\label{lesson:PT-C242-claro-clara}



\begin{quote}
Before the new one: say the Portuguese for ripe; mature, then the Portuguese for dark.
\end{quote}

\subsection*{You'll want to know: claro / clara}
\textbf{claro / clara} — \textquotedblleft{}clear, pale (of a colour)\textquotedblright{}.

\begin{grammarlens}[title={What you've built}]
One more word for this chapter.
\end{grammarlens}

\subsection*{Your turn}
\begin{itemize}
\raggedright
  \item \emph{Say it:} \emph{claro / clara}
  \item \emph{Say it:} \emph{claro / clara}, once more
  \item \emph{Say it:} say \emph{claro / clara}
  \item \emph{Recall:} say the Portuguese for ripe; mature, then the Portuguese for dark, then say \emph{claro / clara} again
\end{itemize}

\begin{tcolorbox}[breakable,colback=teal!4,colframe=teal!35!black,title={Before you move on}]
What does claro / clara mean? (\textquotedblleft{}Clear, pale (of a colour)\textquotedblright{}.) Say it once more.
\end{tcolorbox}
\section[distante]{distante — distant}

\label{lesson:PT-C242-distante}



\begin{quote}
Before the new one: say the Portuguese for dark, then the Portuguese for clear.
\end{quote}

\subsection*{You'll want to know: distante}
\textbf{distante} — \textquotedblleft{}distant\textquotedblright{}.

\begin{grammarlens}[title={What you've built}]
One more word for this chapter.
\end{grammarlens}

\subsection*{Your turn}
\begin{itemize}
\raggedright
  \item \emph{Say it:} \emph{distante}
  \item \emph{Say it:} \emph{distante}, once more
  \item \emph{Say it:} say \emph{distante}
  \item \emph{Recall:} say the Portuguese for dark, then the Portuguese for clear, then say \emph{distante} again
\end{itemize}

\begin{tcolorbox}[breakable,colback=teal!4,colframe=teal!35!black,title={Before you move on}]
What does distante mean? (\textquotedblleft{}Distant\textquotedblright{}.) Say it once more.
\end{tcolorbox}
